% Original applied explanation; see Chapter04/SOURCES.md.
\subsubsection{Notice লিংক ভুল ফোল্ডারে যায় কেন?}
ধরা যাক পেজটি \texttt{/students/index.html}-এ আছে। সেখানে \texttt{href="notice.html"} দিলে সাধারণ অবস্থায় \texttt{/students/notice.html} খোঁজা হবে; \texttt{../notice.html} দিলে এক স্তর ওপরের \texttt{/notice.html} খোঁজা হবে। তাই আপেক্ষিক পথ ঠিক কি না, তা বর্তমান পেজের ভিত্তি ঠিকানা ধরে বিচার করবে।

দর্শককে গন্তব্য বোঝাতে \textbf{ভর্তির বিজ্ঞপ্তি পড়ুন} ধরনের লিংক-লেখা ব্যবহার করবে। একই সাইটে সব লিংকে \textbf{এখানে ক্লিক করুন} লিখলে কোন লিংক কোন তথ্য দেবে তা বোঝা কঠিন হয়।
